Let $X$ be a smooth projective variety of dimension $n$ defined over $k = \mathbf{F}_q$. Let $Z(t)$ be the zeta function of $X$. Then:

\textbf{1.1. Rationality.} $Z(t)$ is a rational function of $t$, i.e., a quotient of polynomials with rational coefficients.

\textbf{1.2. Functional equation.} Let $E$ be the self-intersection number of the diagonal $\Delta$ of $X \times X$ (which is also the top Chern class of the tangent bundle of $X$ (App. A, Ex. 6.6)). Then $Z(t)$ satisfies

$$Z\left(\frac{1}{q^n t}\right) = \pm q^{nE/2} t^E Z(t).$$

\textbf{1.3. Analogue of the Riemann hypothesis.} It is possible to write

$$Z(t) = \frac{P_1(t) P_3(t) \cdots P_{2n-1}(t)}{P_0(t) P_2(t) \cdots P_{2n}(t)}$$

where $P_0(t) = 1-t$, $P_{2n}(t) = 1-q^n t$, and for $1 \leqslant i \leqslant 2n-1$, $P_i(t)$ is a polynomial with integer coefficients which can be written $P_i(t) = \prod_j (1 - \alpha_{ij} t)$ where the $\alpha_{ij}$ are algebraic integers with $|\alpha_{ij}| = q^{i/2}$.

\textbf{1.4. Betti numbers.} Assuming (1.3), the $i$-th Betti number $B_i = B_i(X)$ is the degree of $P_i(t)$. Then $E = \sum (-1)^i B_i$. If $X$ is obtained from a variety $Y$ over an algebraic number ring by reduction modulo a prime ideal, then $B_i(X)$ equals the $i$-th Betti number of $Y_h = (Y \times_R \mathbf{C})_h$.
